% © 2026 Ing. David Jaroš — CC BY-NC-ND 4.0
%
% This work is licensed under a Creative Commons Attribution-NonCommercial-NoDerivatives
% 4.0 International License (CC BY-NC-ND 4.0).
%
% License History: Earlier drafts (up to v0.3) were released under CC BY 4.0.
% From v0.4 onward, all material is released under CC BY-NC-ND 4.0 to protect
% the integrity of the theoretical work during ongoing academic development.
%
% See LICENSE.md for full license text.
%
% File: research_tracks/coupling_spectrum/rg_prime_checkpoints.tex
%
% Task: map_stable_prime_sectors_to_coupling_constants — Target 3
% Priority: CRITICAL
% Mode: hypothesis_test_no_numerology
%
% Purpose: Test whether stable primes 127 and 137 can be understood as RG
%   checkpoints of the electromagnetic coupling, and whether the remaining
%   stable primes {139, 151, 157} align with any coarse-grained RG structure.
%
% Hard rules:
%   - Do NOT use any prime to derive the coupling.
%   - Do NOT claim derivation.
%   - Matches are hypotheses; flag all as OBSERVED CONSISTENCY.

% UBT-AI-PROVENANCE-BEGIN
% schema: ubt-ai-provenance/v1
% tier: C_working
% ai_assistance: disclosed
% human_review: risk-based
% editorial_responsibility: Ing. David Jaroš
% policy: ../../AI_PROVENANCE.md
% notice: Working material; exhaustive human review is not claimed.
% UBT-AI-PROVENANCE-END

\documentclass[12pt,a4paper]{article}
\usepackage[a4paper, margin=2.5cm]{geometry}
\usepackage{amsmath, amssymb, amsthm}
\usepackage{mathtools}
\usepackage{hyperref}
\usepackage{xcolor}
\usepackage{booktabs}
\usepackage{longtable}
\usepackage{mdframed}

\newtheorem{proposition}{Proposition}[section]
\theoremstyle{definition}
\newtheorem{definition}[proposition]{Definition}
\newtheorem{hypothesis}[proposition]{Hypothesis}
\theoremstyle{remark}
\newtheorem{remark}[proposition]{Remark}

\newmdenv[backgroundcolor=yellow!12, linecolor=orange!70, linewidth=1.5pt]{gapbox}
\newmdenv[backgroundcolor=green!8,  linecolor=green!55!black, linewidth=1.5pt]{resultbox}
\newmdenv[backgroundcolor=red!7,    linecolor=red!55,   linewidth=1.5pt]{warnbox}
\newmdenv[backgroundcolor=blue!5,   linecolor=blue!45,  linewidth=1.5pt]{obsbox}

\newcommand{\cond}{\textcolor{blue!65!black}{\textbf{[COND]}}}
\newcommand{\obs}{\textcolor{blue!45!black}{\textbf{[OBS-CONSISTENCY]}}}
\newcommand{\noev}{\textcolor{red!65!black}{\textbf{[NO-EVIDENCE]}}}
\newcommand{\hyp}{\textcolor{orange!75!black}{\textbf{[HYPOTHESIS]}}}
\newcommand{\der}{\textcolor{green!55!black}{\textbf{[DERIVED]}}}

\title{\textbf{RG Prime Checkpoints}\\[4pt]
       \large Stable Primes as Checkpoints on the Electromagnetic RG Trajectory\\[2pt]
       \normalsize Task: \texttt{map\_stable\_prime\_sectors\_to\_coupling\_constants} — Target 3}
\author{Ing.\ David Jaro\v{s}\\
  \small \texttt{research\_tracks/coupling\_spectrum/}}
\date{May 2026}

\begin{document}
\maketitle

\begin{abstract}
We test whether the stable-prime set $\mathcal{S} = \{2, 127, 137, 139, 151, 157\}$
can be interpreted as coarse-grained checkpoints on the electromagnetic RG
trajectory.
First, the QED/SM running of $\alpha_{\rm em}^{-1}$ is derived from the
one-loop beta function without using any prime as input.
The resulting trajectory is then compared to the stable-prime set.
Two primes (127 and 137) fall within the physical running range of
$\alpha_{\rm em}^{-1}$.
The RG trajectory passes continuously from $\alpha_{\rm em}^{-1}(0) \approx 137.036$
to $\alpha_{\rm em}^{-1}(M_Z) \approx 127.9$ over approximately 11 orders of
magnitude in energy, and this flow is consistent with — but does not derive —
the presence of 127 and 137 in $\mathcal{S}$.
Primes 139, 151, and 157 do not align with any known RG checkpoint.
\textbf{Verdict}: OBSERVED\_CONSISTENCY for the 127/137 pair;
NO\_EVIDENCE for the remaining primes.
\end{abstract}

\tableofcontents
\clearpage

% -----------------------------------------------------------------------
\section{Setup and Hard Rules}

\begin{warnbox}
\textbf{Hard rules (enforced throughout)}:
\begin{itemize}
\item No stable prime is used to derive any coupling constant.
\item No coupling constant is adjusted to match any prime.
\item All proximity observations are classified as hypotheses or observed consistency,
      not derivations.
\item The phrase ``alpha = 1/137 because 137 is a stable prime'' is explicitly prohibited.
\end{itemize}
\end{warnbox}

The stable-prime set from the UBT potential $V(q; B(p)) = q^2 - B(p)\,q\ln q$
with $B(p) = (p+1)/3$ is:
\[
  \mathcal{S} = \{2,\; 127,\; 137,\; 139,\; 151,\; 157\}.
\]
This set is derived purely from the mathematical structure of $V$ and contains
no physical input.

% -----------------------------------------------------------------------
\section{One-Loop RG Equation for the Electromagnetic Coupling}

\subsection{Beta function}

In QED with $N_f$ Dirac fermions of charges $Q_r$ in units of $e$, the
one-loop beta function for $\alpha_{\rm em} = e^2/(4\pi)$ is:
\begin{equation}
  \mu \frac{d\alpha_{\rm em}}{d\mu}
  = \frac{b_{\rm em}}{2\pi}\,\alpha_{\rm em}^2 + O(\alpha_{\rm em}^3),
  \label{eq:beta_alpha}
\end{equation}
where
\begin{equation}
  b_{\rm em} = \frac{4}{3}\sum_{\text{Dirac}} Q_r^2
             + \frac{1}{3}\sum_{\text{scalar}} Q_r^2.
  \label{eq:bem}
\end{equation}
Rewriting in terms of the inverse coupling:
\begin{equation}
  \frac{d\alpha_{\rm em}^{-1}}{d\ln\mu} = -\frac{b_{\rm em}}{2\pi}.
  \label{eq:beta_inv}
\end{equation}
Since $b_{\rm em} > 0$ for any non-trivial charged matter content,
$\alpha_{\rm em}^{-1}$ \emph{decreases} monotonically with increasing $\mu$.

\begin{resultbox}
\der\ The sign of the running is fixed: $\alpha_{\rm em}^{-1}$ is a
decreasing function of the energy scale $\mu$.  This holds for any positively
charged matter field and requires no experimental input.
\end{resultbox}

\subsection{Integrated RG trajectory}

Integrating \eqref{eq:beta_inv} from reference scale $\mu_0$ to $\mu$:
\begin{equation}
  \alpha_{\rm em}^{-1}(\mu)
  = \alpha_{\rm em}^{-1}(\mu_0)
  - \frac{b_{\rm em}}{2\pi}\,\ln\frac{\mu}{\mu_0}
  + O(\alpha_{\rm em}).
  \label{eq:rg_solution}
\end{equation}
In the full SM, $b_{\rm em}$ varies with the energy scale as successive
particle thresholds are crossed.  Below each threshold $m_f$, the
corresponding fermion decouples and its contribution to $b_{\rm em}$ switches
off.

\subsection{Effective beta coefficient in the SM}

The SM particle content (active at scale $\mu$) contributes to
$b_{\rm em}$ as follows:

\begin{center}
\begin{tabular}{lll}
\toprule
Particle & $Q_r$ & Contribution to $b_{\rm em}$ \\
\midrule
Electron $e$ & $-1$ & $4/3$ \\
Muon $\mu$ & $-1$ & $4/3$ (above $m_\mu$) \\
Tau $\tau$ & $-1$ & $4/3$ (above $m_\tau$) \\
$u,c,t$ quarks (3 colours) & $+2/3$ & $3 \times (4/3)(4/9) = 16/9$ each \\
$d,s,b$ quarks (3 colours) & $-1/3$ & $3 \times (4/3)(1/9) = 4/9$ each \\
\bottomrule
\end{tabular}
\end{center}

Total at $\mu \gg m_t$ (all SM fermions active):
\[
  b_{\rm em}^{\rm SM}
  = \underbrace{3\cdot\tfrac{4}{3}}_{\text{leptons}}
  + \underbrace{3\cdot\tfrac{16}{9}}_{\text{up quarks}}
  + \underbrace{3\cdot\tfrac{4}{9}}_{\text{down quarks}}
  = 4 + \tfrac{48}{9} + \tfrac{12}{9}
  = 4 + \tfrac{20}{3}
  = \tfrac{32}{3} \approx 10.67.
\]

% -----------------------------------------------------------------------
\section{Trajectory Values and Comparison with Stable Primes}

\subsection{Trajectory from first principles (qualitative)}

Using \eqref{eq:rg_solution} with $\mu_0 = m_e$ and the one-electron
approximation $b_{\rm em} = 4/3$:
\[
  \alpha_{\rm em}^{-1}(M_Z)
  \approx \alpha_{\rm em}^{-1}(m_e)
  - \frac{4/3}{2\pi}\,\ln\frac{91\,200}{0.511}
  \approx 137.036 - \frac{4/3}{2\pi} \times 12.09
  \approx 137.036 - 2.57
  \approx 134.5.
\]
This underestimates the full SM running because only the electron is included.
The full SM value (including all thresholds and hadronic contributions) gives:
\[
  \alpha_{\rm em}^{-1}(M_Z) \approx 127.9.
\]
The difference (137.036 − 127.9 = 9.14) is accounted for by muon, tau,
quark, and hadronic threshold contributions.

\subsection{RG trajectory table}

\begin{center}
\begin{tabular}{llll}
\toprule
Scale $\mu$ & $\alpha_{\rm em}^{-1}(\mu)$ & Stable prime in $\mathcal{S}$? & Distance \\
\midrule
$0$ (Thomson) & $137.036$ & $137 \in \mathcal{S}$ & $0.036$ \\
$m_e = 0.511$ MeV & $137.036$ & $137 \in \mathcal{S}$ & $0.036$ \\
$m_\mu = 105.66$ MeV & $\approx 136.0$ & none & — \\
$m_\tau = 1.777$ GeV & $\approx 133.5$ & none & — \\
$M_Z = 91.19$ GeV & $127.9$ & $127 \in \mathcal{S}$ & $0.9$ \\
\bottomrule
\end{tabular}
\end{center}

\begin{obsbox}
\obs\ The RG trajectory of $\alpha_{\rm em}^{-1}$ passes through the numerical
neighbourhood of $p = 137$ at low energy and $p = 127$ at the $Z$ mass.
Neither value is reached exactly; the distances are 0.036 and 0.9 respectively.
\end{obsbox}

\subsection{Can 127 and 137 be RG checkpoints?}

\begin{hypothesis}[127–137 as RG checkpoints]
The primes 127 and 137 in $\mathcal{S}$ correspond to approximate checkpoints
on the electromagnetic running:
\begin{itemize}
\item $p = 137$: $\alpha_{\rm em}^{-1} \approx 137$ at $\mu \lesssim m_e$ (0.026\% error).
\item $p = 127$: $\alpha_{\rm em}^{-1} \approx 128$ at $\mu = M_Z$ (0.70\% error).
\end{itemize}
\end{hypothesis}

\textbf{Assessment}:
\begin{itemize}
\item The trajectory \emph{is} continuous and \emph{does} pass from $\approx 137$ to
      $\approx 128$ over the known energy range.
\item The stable primes 127 and 137 are consecutive integers separated by 10,
      closely bracketing the physical running window of $9.14$ units.
\item However, the precise physical values are $\alpha_{\rm em}^{-1}(0) = 137.036$
      (not $137$) and $\alpha_{\rm em}^{-1}(M_Z) = 127.9$ (not $127$).
\item These offsets (0.036 and 0.9 respectively) are non-zero and cannot be
      accounted for by prime-stability alone.
\item There is no known UBT mechanism that selects integer-valued $\alpha^{-1}$ at
      specific scales.
\end{itemize}

\textbf{Conclusion for this pair}: The electromagnetic coupling trajectory
passes through the neighbourhood of stable primes 127 and 137, giving
OBSERVED\_CONSISTENCY.  This is not derivation; it is a proximity observation.

% -----------------------------------------------------------------------
\section{Non-Matching Stable Primes}

\subsection{Prime p = 139}

The value $\alpha_{\rm em}^{-1}$ at no standard energy scale equals 139.
The trajectory passes through $\alpha_{\rm em}^{-1} = 139$ at some scale
below the electron mass (in the deep infrared, effectively at $\mu \to 0$
where $\alpha_{\rm em}$ is constant).  The Thomson-limit value is 137.036,
not 139.

Distance to nearest coupling: 1.964 from $\alpha_{\rm em}^{-1}(0)$ (1.43\%).
This exceeds experimental precision by a factor $\sim 10^4$ (the value
137.036 is known to ppm precision).

\noev\ No physical RG checkpoint at $\alpha_{\rm em}^{-1} = 139$.

\subsection{Prime p = 151}

$\alpha_{\rm em}^{-1} = 151$ would correspond to a scale below the electron
mass in the direction of negative $\mu$ or equivalently below the Landau pole
of QED — unphysical territory.  No SM coupling has $\alpha^{-1} \approx 151$
at any accessible energy scale.

\noev\ No physical RG checkpoint at $\alpha_{\rm em}^{-1} = 151$ or equivalent.

\subsection{Prime p = 157}

Similarly, no SM coupling inverse is near 157 at any physical scale.

\noev\ No physical RG checkpoint at $\alpha_{\rm em}^{-1} = 157$ or equivalent.

% -----------------------------------------------------------------------
\section{Coarse-Grained RG Plateaus: Can the Whole Set Be Interpreted?}

A stronger interpretation would be that the full set $\mathcal{S}$ corresponds
to coarse-grained plateaus or resonance windows in a \emph{generalised} RG
flow — perhaps in a UBT multi-sector extension with running couplings in
additional gauge groups or Kaluza--Klein towers.

Requirements for this interpretation:
\begin{enumerate}
\item A UBT mechanism that produces multiple effective $\alpha_i^{-1}$ values
      in the range 127–157.
\item A selection rule identifying which values in that range are stable
      (analogous to the prime stability condition in the UBT potential).
\item An energy scale assignment for each element of $\mathcal{S}$.
\end{enumerate}

\textbf{Current status}:
\begin{itemize}
\item No UBT mechanism is currently known that produces effective
      $\alpha_i^{-1} \in \{139, 151, 157\}$.
\item The weak SU(2) and hypercharge U(1) couplings at $M_Z$ give
      $\alpha_2^{-1} \approx 29.6$ and $\alpha_1^{-1} \approx 58.7$,
      which are far from $\{139, 151, 157\}$.
\item The strong coupling $\alpha_3^{-1}(M_Z) \approx 8.47$ is also not near
      any of these primes.
\end{itemize}

\begin{gapbox}
\textbf{Open gap}: The stable primes 139, 151, 157 have no identified
correspondence in the Standard Model coupling spectrum or in the UBT
multi-sector extension.  Whether they can be connected to couplings in
compactified sectors, Hecke towers, or Kaluza--Klein spectra is an open
question requiring further derivation.
\end{gapbox}

% -----------------------------------------------------------------------
\section{Verdict}

\begin{center}
\begin{tabular}{lll}
\toprule
Prime & Physical coupling & Classification \\
\midrule
$p = 2$ & None (structural) & STRUCTURAL ELEMENT \\
$p = 127$ & $\alpha_{\rm em}^{-1}(M_Z) \approx 127.9$ & OBS-CONSISTENCY \\
$p = 137$ & $\alpha_{\rm em}^{-1}(0) \approx 137.036$ & OBS-CONSISTENCY \\
$p = 139$ & None at any SM scale & NO-EVIDENCE \\
$p = 151$ & None at any SM scale & NO-EVIDENCE \\
$p = 157$ & None at any SM scale & NO-EVIDENCE \\
\bottomrule
\end{tabular}
\end{center}

\textbf{Combined RG verdict}:
The electromagnetic coupling runs continuously from $\sim 137$ to $\sim 128$
over the energy range $[m_e, M_Z]$.  The stable-prime set contains both
endpoints (approximately).  This is OBSERVED\_CONSISTENCY, not a derivation.
The remaining three primes (139, 151, 157) are currently unmapped.

\begin{resultbox}
\textbf{Mandatory final sentence}:\\
Stable primes do not currently map to Standard Model couplings beyond
alpha-like coincidences.
\end{resultbox}

\end{document}
